% Generated draft: requires mathematical and source review.
\documentclass[11pt]{article}
\usepackage[margin=25mm]{geometry}
\usepackage{fontspec}
\setmainfont{Latin Modern Roman}
\usepackage{amsmath,amssymb,mathtools,unicode-math}
\setmathfont{Latin Modern Math}
\usepackage{enumitem,needspace,xurl,hyperref}
\setlength{\parindent}{0pt}
\setlength{\parskip}{5pt}
\setlength{\emergencystretch}{3em}
\newcommand{\N}{\mathbb{N}}
\newcommand{\Z}{\mathbb{Z}}
\newcommand{\Q}{\mathbb{Q}}
\newcommand{\R}{\mathbb{R}}
\newcommand{\C}{\mathbb{C}}
\newcommand{\F}{\mathbb{F}}
\newcommand{\A}{\mathbb{A}}
\newcommand{\PP}{\mathbb{P}}
\newcommand{\E}{\mathbb{E}}
\newcommand{\eps}{\varepsilon}
\newcommand{\dd}{\,\mathrm{d}}
\begin{document}
\section*{UnsolvedMath ID: 30000881}
\textbf{Source problem code:} catalogue entry
\subsection{Definitions and mathematical statement}
The record combines three separate questions; their groups and coefficient fields need not be the same.

First, let \(G\) be a finite \(p\)-group. The source assumes that \(G\), embedded as group elements in \((kG)^\times\), is a ``complement factor'' in that unit group, and asks
\[
 kG\cong kH\text{ as }k\text{-algebras}\quad\Longrightarrow\quad G\cong H\ ?
\]
The field \(k\), its characteristic, the group class for \(H\), and the complement-factor condition are not defined. A normal-complement condition would mean \((kG)^\times=NG\), \(N\triangleleft(kG)^\times\), \(N\cap G=1\); a direct-factor condition would additionally require commuting factors. Neither interpretation is silently selected, nor is the full unit group replaced by its augmentation-one subgroup.

Second, take \(G=\operatorname{GL}_n(k)\), where \(k\) is a finite field and \(n\geq1\). For \(u=\sum_{g\in G}a_g g\in\Z G\), put
\[
 \varepsilon(u)=\sum_ga_g,\qquad
 V(\Z G)=\{u\in(\Z G)^\times:\varepsilon(u)=1\},\qquad
 \varepsilon_C(u)=\sum_{g\in C}a_g
\]
for each conjugacy class \(C\) of \(G\). A torsion unit has finite multiplicative order. The question is
\[
 \forall u\in V(\Z G)\text{ of finite order},\qquad
 \varepsilon_{\pi(u)}(u)\ne0\ ?
\]
The source does not define \(\pi\), or whether its subscript denotes an element's conjugacy class. A natural candidate is the map induced by \(\Z G\to\operatorname{Mat}_n(k)\), sending each group element to its matrix and integer coefficients to \(k\); it sends units to \(\operatorname{GL}_n(k)\). If that is the intended map, the subscript means the conjugacy class of this matrix. The candidate is identified only to expose the missing convention, not asserted as a recovered definition.

Third, let \((u_p)_{p\in\mathcal P}\) be the nonempty family of torsion units constructed in the report from central primitive idempotents of localized cyclotomic group algebras, viewed in \(\C G\). The construction, local rings, and index set \(\mathcal P\) are absent. The question is whether
\[
 [\ \forall p,q\in\mathcal P\ \exists x_{p,q}\in(\C G)^\times:
 x_{p,q}^{-1}u_px_{p,q}=u_q\ ]
\]
\[
 \Longrightarrow\quad
 [\ \exists g\in G\ \forall p\in\mathcal P\ \exists y_p\in(\C G)^\times:
 y_p^{-1}u_py_p=g\ ].
\]
By mutual conjugacy, being conjugate to a group element for one index would imply it for every index, with the same element \(g\) after appropriate conjugators. Conjugacy is explicitly over \(\C G\); it is not silently changed to \(\Q G\) or integral conjugacy.
\subsection{Short English statement}
The three questions concern recovering a finite group from its group algebra under an unspecified complement condition, nonvanishing of a particular partial augmentation for general linear groups, and whether a specially constructed mutually conjugate family of units is conjugate to group elements. Their missing maps and construction hypotheses prevent a complete unified theorem statement.
\subsection{Sources}
\begin{itemize}
\item Local UnsolvedMath registry, permanent ID \(30000881\), exact source code \texttt{catalogue entry}, proposed year 2007; no individual proposer is supplied.
\item \emph{Mini-Workshop: Arithmetik von Gruppenringen (2008)}, Oberwolfach Reports, DOI 10.4171/owr/2007/55, pp. 3209--3240. \url{https://doi.org/10.4171/owr/2007/55}.
\item Supplied reference labelled ``rational conjugacy for nonsolvable groups'': \url{https://arxiv.org/abs/1305.7419}.
\item Supplied reference labelled ``Zassenhaus counterexamples'': \url{https://doi.org/10.1016/j.aim.2018.10.004}.
\item The supplied \texttt{corrected\_verified} note says all three questions were retained. The assessment checked 2026-08-21 explicitly says the record conflates modular algebra-isomorphism, partial-augmentation, and rational-conjugacy problems, reports substantial partial results, and labels the record partially solved. This is dated supplied triage; the counterexample reference and precise scope of remaining general questions require separate verification.
\end{itemize}
\end{document}
